\documentclass[11pt]{article}
\usepackage[T1]{fontenc}
\usepackage{lmodern}
\usepackage[margin=0.92in]{geometry}
\usepackage{amsmath,amssymb,amsthm,booktabs,microtype}
\usepackage[colorlinks=true,linkcolor=blue,citecolor=blue,urlcolor=blue]{hyperref}
\hypersetup{pdftitle={Integrated geometric control of shell-volume response},pdfauthor={Ricardo Maldonado},pdfsubject={Working draft for author and specialist review}}
\pdfinfoomitdate=1
\pdftrailerid{}
\pdfsuppressptexinfo=15
\newcommand{\eps}{\varepsilon}
\newcommand{\dd}{\mathrm d}
\newcommand{\Lie}{\mathcal L}
\newcommand{\norm}[1]{\lVert #1\rVert}
\newcommand{\cV}{\mathcal V}
\newcommand{\HH}{\mathsf H}
\newcommand{\ip}[2]{\langle #1,#2\rangle_{\HH}}
\newcommand{\PN}{P_N}
\newtheorem{theorem}{Theorem}
\newtheorem{lemma}{Lemma}
\newtheorem{proposition}{Proposition}
\title{Integrated geometric control of shell-volume response}
\author{Ricardo Maldonado\\\small Independent Researcher}
\date{4 November 2026}

\begin{document}
\maketitle
\begin{center}\small Working draft for author and specialist review\end{center}
\begin{abstract}
A derivative-supremum condition can exclude concentrated smooth initial
data even when their integrated geometric size remains bounded. We derive
a direct $H^1$-to-observable estimate for the retained five-channel
Einstein--scalar shell mean problem. On a physically supplied initial
Cauchy section, the profile $D$ is the relative difference of second
scalar-level three-area variations, with fixed first perturbation and
background scalar-label calibration. A tilted wave energy with a positive
boundary value term controls the scalar history despite a mixed
Dirichlet and nonselfadjoint Robin boundary block. The complete initial
map gives $E(0)\le R^2/(5\times10^6)$, and the boundary history obeys
$|m|\le R/300$, where $R=\norm{D}_{H^1([1,5/4])}$. The specified closed
class $R\le2\times10^{31}$ retains strictly negative volume and expansion
second coefficients at one proper time. An explicit sequence in the
actual compact preparation image has bounded $H^1$ norm and unbounded
old derivative-supremum norm, so this criterion covers data beyond that
earlier sufficient criterion. It does not imply containment of the two
budget balls or response instability. All coefficient, exact-family and
linear-development premises remain explicit and conditional. The estimating
energy is not an identified physical energy or preparation law; neither a
physical class selection nor a uniform nonlinear parameter interval follows.
\end{abstract}

\section{Model, fixed data, and the new estimate}

The sign of a second response depends on both the common quadratic first
source and the homogeneous second preparation. The retained shell model
admits an explicit pair with opposite volume and expansion coefficient
signs, at the same anchored proper time, while its first packet is fixed
\cite{preparation}. A subsequent geometric class, based on the supremum of
an initial scalar-level area profile and its derivative, certifies negative
signs throughout a specified budget \cite{intrinsic}. The question here is
whether integrated control of that same profile can give a direct response
bound, even when the earlier derivative-supremum route fails.

The bulk normalization and one-copy outward shell equations are
\begin{align}
 S_{\rm bulk}&=\int\bigl(R_g/2-|\dd\phi|_g^2/2-\cV(\phi)\bigr)\dd\mu_g,
 \nonumber\\
 \operatorname{Ric}_{AB}&=\phi_A\phi_B+\tfrac23\cV g_{AB},
 &\Box_g\phi&=\cV',\nonumber\\
 K^{\rm out}_{ab}&=\tfrac16\sigma(\phi)\gamma_{ab},
 &n_{\rm out}\phi&=-\tfrac12\sigma'(\phi).
 \label{eq:model}
\end{align}
Two identical interior copies use their respective outward normals. Fix
\begin{gather}
 W=1-\phi+\phi^3/3,\qquad
 \cV=(W')^2/2-2W^2/3,\qquad
 \sigma=2W+10^{-3}(1+c_\sigma\phi),\nonumber\\
 c_\sigma=96714321425692/161839258930731.
 \label{eq:potential}
\end{gather}
The selected regular background and coordinates are
\begin{equation}
 g_0=\dd y^2+\rho(y)^2[-\dd t^2+\cosh^2t\,\kappa],\qquad
 a=\rho_y/\rho>0,\quad p=\phi_{0,y}\ne0,\quad
 \dd y=-\ell\rho\,\dd r,\quad x=t/\ell,
 \label{eq:background}
\end{equation}
where $\kappa$ is the unit round three-sphere metric, $\ell=10^{-6}$,
the packet is $1\le r\le2$, and the exact selected shell radius is $\rho_b$.
The target is
\begin{equation}
 T=5/4,\qquad \Delta=T-1=1/4,\qquad \tau_*=\rho_b\ell T.
 \label{eq:target}
\end{equation}
The exact $\rho_b$ fixes the clock; its enclosure is not a replacement
value. All numerical quantities below are in the model's fixed parameter
and coefficient normalization, with no SI calibration.

We compare the same background, first canonical quadrupole seed, complete
quadratic constraint and evolution sources, initial shell collar, canonical
momentum and both scalar canonical second sectors, and second
transverse-traceless initial pair. Parameter derivatives use
\begin{equation}
 x_\eps=x_0+\eps v+\tfrac12\eps^2w+o(\eps^2).
 \label{eq:normalization}
\end{equation}
The common smooth flat seed has $Y=n_1^2-n_2^2$,
$\langle Y\rangle=0$, $\langle Y^2\rangle=1/6$ and
$\Delta_\kappa Y=-8Y$; its exact normalization and the fixed base
preparation are those specified in \cite{preparation,intrinsic}.
The retained base response includes the complete first-field bulk and
shell forcing, clock corrections and measure corrections.

We assume the selected smooth background and its authenticated continuous
boxes; the equivariant weighted constraint inverse and nonlinear
implicit-function hypotheses; the conditional exact-family fixed-shell
development, subsidiary and boundary uniqueness; and identification of the
common first derivative. Each fixed smooth preparation has $C^2$ parameter
dependence in the retained lower $C^4$ topology on a compact neighborhood.
An anchored intrinsic normal tube reaches the target after a possibly
family-dependent shrinkage in $\eps$. The exact five-mean projection,
homogeneous difference equation, boundary rows and continuous operator
caps in Section~\ref{sec:energy} are retained conditional premises.
In particular the numerical caps are declared by the earlier
\texttt{EXPLICIT\_PREMISES.json} record \cite{records}; they are not
reproved by the new arithmetic. The new result is an a priori estimate
for this stipulated coupled problem, not a new nonlinear existence theorem.

\section{Geometric profile and complete initial map}

On the supplied physical initial section $\Sigma$, let $h$ be smooth,
compact in the packet, flat at both endpoints, and satisfy
\begin{equation}
 \int_{y_L}^{y_U}\rho^2h\,\dd y=0,\qquad
 \chi(y)=\int_{y_L}^y\rho^2h\,\dd s,\qquad
 h=\rho^{-2}\chi_y,\quad y_L=y(2),\quad y_U=y(1).
 \label{eq:admissible}
\end{equation}
The complete added raw canonical second datum is
\begin{equation}
 (\Delta\gamma_2,\Delta\pi_2,\Delta\phi_2,\Delta\varpi_2)
 =\bigl(h\gamma_0+\Lie_{\zeta\partial_y}\gamma_0,0,0,0\bigr),
 \qquad
 \zeta=\frac{3(h_y-ah)}{2p^2}+\frac{9\chi}{2\rho^4p^2}.
 \label{eq:raw}
\end{equation}
Its exact homogeneous constraint and zero-moment construction, and its
conditional exact completion with the same full correction second
derivative, are retained from \cite{preparation}. No quadratic source is
changed by this addition.

For $\varphi=\phi_0(y)$, use the scalar-level component in the annular
collar that continues the background sphere. Its three-area is
$A_\eps(\varphi)$ and $A_0(\phi_0(y))=2\pi^2\rho^3$. At fixed scalar
label define
\begin{equation}
 D=\frac{\partial_\eps^2A_\eps^h|_0-
                    \partial_\eps^2A_\eps^{\rm base}|_0}{A_0},
 \qquad
 D_r=(-\ell\rho p)\partial_\varphi D.
 \label{eq:geometry}
\end{equation}
This derivative uses the fixed background scalar-label map, not an arbitrary
coordinate of the perturbed metric. In a slightly enlarged collar, $p\ne0$
gives regular continued level graphs. The retained $C^2$-into-$C^4$
dependence gives $C^2$-into-$C^1$ area dependence after the two spatial
orders of loss described in \cite{intrinsic}. Thus $D$ and $D_r$ exist.
There is no assumption of exact rotational symmetry of the perturbed data.

The two first level displacements agree. Their second displacement
difference is $-\Delta\phi_2/p$, and the common first-field determinant
quadratics cancel after subtraction. For the angular metric mean $J$ and
scalar mean $w$ this gives
\begin{equation}
 D=\tfrac32J-3aw/p.
 \label{eq:D}
\end{equation}
The raw representative has $J=h+2a\zeta,w=0$; after the pure second
conditioning gauge it has $J=h,w=-p\zeta$. Both give the same $D$.
This is the standard metric--scalar displacement cancellation on the
given section. Simultaneous diffeomorphisms transport the section, fields
and level component; choosing a different physical section is a different
observable.

Set $Z=\chi/(\ell\rho^3)$, $b=\ell\rho a$,
$g=\ell\rho p^2/(3a)>0$, and $\beta=p/(2a)$. Then
\begin{equation}
 Z_r=3bZ-h,\qquad
 h_r=3\ell^2Z+(-b+g)h-\tfrac23gD,\qquad (Z,h)(1)=(0,0).
 \label{eq:reconstruct}
\end{equation}
This retained reconstruction determines the restricted preparation on
$[1,T]$ from its profile there; whole-packet injectivity needs the full
profile on $[1,2]$. It is not injectivity on unrestricted Cauchy data or
surjectivity onto arbitrary profiles.

In $(K,L,N,J,w)$ order the complete conditioned difference initial map is
\begin{gather}
 U_0=(h,0,0,h,w),\qquad V_0:=U_{0,x}=(0,V_L,0,0,0),\qquad
 P_0:=U_{0,r}=(h_r,0,0,h_r,w_r),\nonumber\\
 w=\beta(h-2D/3),\qquad
 w_r=\beta_r(h-2D/3)+\beta(h_r-2D_r/3),\qquad V_L=h_r+bh.
 \label{eq:initialmap}
\end{gather}
In particular the mixed metric velocity and derivative of the scalar
coefficient are included. Define the integrated profile size
\begin{equation}
 R=\norm{D}_{H^1(I)}
   :=\left(\int_I(D^2+D_r^2)\,\dd r\right)^{1/2},
 \qquad I=[1,T].
 \label{eq:R}
\end{equation}
It applies to the complete increment: $R(\lambda h)=|\lambda|R(h)$.

\section{Tilted five-channel energy and boundary history}\label{sec:energy}

Write $|U|_{\HH}^2=U^\mathsf{T}\HH U$ for
$\HH=\operatorname{diag}(1,2,1,3,1)$. All operator norms in this section
are the induced norms of this inner product. Retain the exact homogeneous
equation
\begin{equation}
 U_{xx}-U_{rr}=\mathcal A U_x+\mathcal B U_r+\mathcal D U,\qquad
 \norm{\mathcal A},\norm{\mathcal B}\le10^{-4},\quad
 \norm{\mathcal D}\le10^{-6}.
 \label{eq:PDE}
\end{equation}
At $r=0$, the $L$ component is Dirichlet and the four complementary
components obey the Robin row
\begin{equation}
 U_L=0,\qquad \PN U_r=M\PN U,\qquad \norm{M}_{\HH}\le1/2000,
 \label{eq:boundary}
\end{equation}
where $\PN$ projects off the $L$ component. The block $M$ need not be
self-adjoint or constant. The derivative in the Robin row is the
increasing-$r$ derivative in \eqref{eq:background}, with that sign already
included in $M$. The retained cap domain is
$0\le x\le3/2$, $\max(0,x-5/4)\le r\le7/4-x$.
The cone
$\mathcal K_T=\{0\le x\le T,\ 0\le r\le T-x\}$ lies in the retained
coefficient domain; the solution is exactly zero for $x+r\le1$.

Let $V=U_x$, $P=U_r$, $k=\alpha=1/2$, and define
\begin{align}
 e&=\tfrac12\bigl(|V|_{\HH}^2+|P|_{\HH}^2+
                     2k\ip{V}{P}+|U|_{\HH}^2\bigr),\nonumber\\
 \mathcal E(x)&=\int_0^{T-x}e\,\dd r+
                         \tfrac{\alpha}{2}|\PN U(x,0)|_{\HH}^2.
 \label{eq:energy}
\end{align}
This is a coercive estimating norm for the linear difference. It is not
identified with a positive conserved physical energy of the full shell.
Pointwise,
\begin{equation}
 e\ge\tfrac{1-k}{2}(|V|_{\HH}^2+|P|_{\HH}^2)+\tfrac12|U|_{\HH}^2.
 \label{eq:coercive}
\end{equation}

\begin{lemma}[Continuous energy and trace bound]\label{lem:energy}
Under \eqref{eq:PDE}--\eqref{eq:boundary}, the retained smooth solution
satisfies
\begin{equation}
 \mathcal E(x)\le e^{Cx}\mathcal E(0),\qquad
 C=201/100,\qquad e^{CT}<13.
 \label{eq:growth}
\end{equation}
If $\mathcal E(0)\le R^2/(5\times10^6)$, its scalar shell history
$m(x)=U_w(x,0)$ satisfies
\begin{equation}
 m(x)=0\quad(x\le1),\qquad |m(x)|\le R/300\quad(1\le x\le T).
 \label{eq:trace}
\end{equation}
\end{lemma}

\begin{proof}
Put $F=\mathcal A V+\mathcal B P+\mathcal D U$. Direct differentiation
gives the local identity
\begin{equation}
 e_x=f_r+\ip{V+kP}{F}+\ip{U}{V},\qquad
 f=\ip{V}{P}+\tfrac k2(|V|_{\HH}^2+|P|_{\HH}^2).
 \label{eq:flux}
\end{equation}
At the shrinking upper endpoint the flux is
\begin{equation}
 f-e=-\tfrac{1-k}{2}|V-P|_{\HH}^2-\tfrac12|U|_{\HH}^2\le0.
 \label{eq:upperflux}
\end{equation}
Thus no boundary data are introduced on that characteristic side.

The Dirichlet row has $U_L=V_L=0$ at the shell and contributes
$-k|P_L|_{\HH}^2/2\le0$. Put $u=\PN U(x,0)$ and $v=\PN V(x,0)$.
The Robin block and the derivative of the boundary value term give
\begin{equation}
 -f_N+\alpha\ip{u}{v}
 =-\tfrac k2|v|_{\HH}^2+\ip{v}{(\alpha I-M)u}
                           -\tfrac k2|Mu|_{\HH}^2.
 \label{eq:Robinflux}
\end{equation}
This identity does not exchange $M$ with its adjoint. Young's inequality
and $\norm{M}\le1/2000$ control the mixed term using only the operator
norm. No derivative of $M$, Robin quadratic potential, or symmetry
assumption is required.
More explicitly, put $s=v+Mu$ and
$\mu_M=1/2000$, $\gamma=\alpha+(1-k)\mu_M$.
The same shell expression becomes
$-k|s|^2/2-(1-k)\ip{s}{Mu}+(1-k)|Mu|^2
+\alpha\ip{u}{s-Mu}$. Hence it is bounded above by
\begin{equation}
 -\tfrac k4|s|_{\HH}^2+
 \left\{\frac{\gamma^2}{k}+(1-k)\mu_M^2+\alpha\mu_M\right\}|u|_{\HH}^2.
 \label{eq:Robinabsorb}
\end{equation}
Relative to the boundary summand of $\mathcal E$, the resulting rate is
$C_{\rm bdry}=2003001/1000000=2.003001<C$.

For completeness, if $a_*=b_*=10^{-4}$ and $d_*=10^{-6}$, the bulk
production is bounded by
$c_V|V|^2+c_P|P|^2+c_U|U|^2$, where
\begin{align}
 c_V&=a_*+(b_*+ka_*)/2+d_*/2+1/2,\nonumber\\
 c_P&=kb_*+(b_*+ka_*)/2+kd_*/2,\qquad
 c_U=d_*/2+kd_*/2+1/2.
 \label{eq:bulkconstants}
\end{align}
Consequently its rate relative to the bulk energy is at most
$\max\{2c_V/(1-k),2c_P/(1-k),2c_U\}
=1000351/500000=2.000702<C$. The shell mixed term in
\eqref{eq:Robinflux} is also absorbed with a rate below $C$ relative to
the positive boundary value term. Combining these estimates gives
$\mathcal E'\le C\mathcal E$ on $0\le x<T$, hence the first inequality
in \eqref{eq:growth}. Continuity gives the endpoint statement.
A positive Taylor sum and geometric remainder give
\begin{equation}
 e^{CT}\le\sum_{j=0}^{24}\frac{(CT)^j}{j!}
       +\frac{(CT)^{25}/25!}{1-CT/26}<13.
 \label{eq:exp}
\end{equation}
Finally $|m|^2\le|\PN U|_{\HH}^2\le2\mathcal E/\alpha$ and
\begin{equation}
 \frac{2\cdot13}{\alpha(5\times10^6)}
       =\frac{13}{1250000}<\frac1{300^2}.
 \label{eq:tracesquare}
\end{equation}
Together with the zero wedge this proves \eqref{eq:trace}.
\end{proof}

\section{Integrated control of all initial components}\label{sec:initial}

The retained continuous background boxes on $I$ give
\begin{gather}
129\le\rho,\rho_b\le131,\quad
3/25000\le-p\le1/7000,\quad111/1000\le a\le14/125,\nonumber\\
|p_y|\le1/4000,\quad|a_y|\le1/10000,\qquad
1/1200\le-F'_b\le1/1000.
\label{eq:caps}
\end{gather}
Here $F'_b$ is the scalar-field derivative at the background shell of
$F=2\cV-\sigma'^2/4+\sigma^2/3$.
Set
\begin{gather}
 b_+=\ell\,131(14/125),\quad
 g_+=\ell\,131(1/7000)^2/[3(111/1000)],\nonumber\\
 H=\Delta(4/3)(2/3)g_+,\qquad
 H_r=3\ell^2H+(b_++g_+)H+(2/3)g_+,\qquad
 V_+=H_r+b_+H.
 \label{eq:ODEcaps}
\end{gather}
The reconstruction matrix has infinity norm at most $20001/20000$.
The retained rational exponential bound on this interval is $4/3$,
so variation of constants gives
$\sup|h|,\sup|Z|\le H\sup|D|$ and
$\sup|h_r|\le H_r\sup|D|$.

On an interval of length $\Delta=1/4$, the elementary one-dimensional
estimate yields
\begin{equation}
 \sup|D|\le2\norm D_2+\tfrac12\norm{D_r}_2
       \le\tfrac{\sqrt{17}}2R\le\tfrac{21}{10}R.
 \label{eq:sobolev}
\end{equation}
Thus define the per-$R$ $L^2$ caps
\begin{align}
 h_2&=(21/20)H,& h_{r,2}&=(21/20)H_r,& v_2^0&=(21/20)V_+,\nonumber\\
 \beta_+&=1/1554,&
 \beta_{r,+}&=\ell\,131\left\{\frac{1/4000}{2(111/1000)}
       +\frac{(1/7000)(1/10000)}{2(111/1000)^2}\right\},\nonumber\\
 w_2&=\beta_+(h_2+2/3),&
 w_{r,2}&=\beta_{r,+}(h_2+2/3)+\beta_+(h_{r,2}+2/3).
 \label{eq:L2caps}
\end{align}
For example $\norm h_2\le h_2R$, while both $\norm D_2$ and
$\norm{D_r}_2$ are bounded by $R$ in the scalar formulas. Equation
\eqref{eq:initialmap} gives
\begin{equation}
 \norm{U_0}_{L^2_{\HH}}\le(2h_2+w_2)R,\quad
 \norm{P_0}_{L^2_{\HH}}\le(2h_{r,2}+w_{r,2})R,\quad
 \norm{V_0}_{L^2_{\HH}}\le2v_2^0R.
 \label{eq:Gramcaps}
\end{equation}
The initial tilt term is exactly zero: $V_0$ occupies only the $L$ slot,
whereas $P_0$ has zero $L$ slot. The shell initial value is zero in the
unchanged collar. Therefore
\begin{equation}
 \mathcal E(0)\le c_0R^2,\quad
 c_0=\tfrac12\bigl[(2h_2+w_2)^2+(2h_{r,2}+w_{r,2})^2
                                  +(2v_2^0)^2\bigr]
                         <\frac1{5000000}.
 \label{eq:E0}
\end{equation}
All these comparisons are rational, including the guard
$17/4<(21/10)^2$. They control all five channels and the complete scalar
pullback, rather than only a free scalar history.

\section{The specified integrated class and the two proper responses}

Use the shell's intrinsic future normal-geodesic foliation, anchored at
the common initial section, and set
\begin{equation}
 V_\eps(\tau)=\int_{S^3}\dd\mu_{h_\eps(\tau)},\qquad
 \Theta_\eps=\partial_\tau\log V_\eps,\qquad
 v_2=\frac{\partial_\eps^2V_\eps|_0}{V_0},\quad
 \Theta_2=\partial_\eps^2\Theta_\eps|_0.
 \label{eq:observables}
\end{equation}
The clock and anchor are transported with a gauge change.
The background has $A(\tau)=\rho_b\cosh(\tau/\rho_b)$,
$V_0=2\pi^2A^3$, and
$\Theta_0=3\rho_b^{-1}\tanh(\tau/\rho_b)>0$ at $\tau_*$.

In the difference from the base, the common first quadratic sources,
proper-clock and measure terms cancel. The retained full coupled
difference gives
\begin{align}
 \nu(T)&:=(v_2-v_2^{\rm base})(T)
   =(\rho_b\ell)^2(F'_b/2)\int_1^T K(T,x)m(x)\,\dd x,\nonumber\\
 \vartheta(T)&:=(\Theta_2-\Theta_2^{\rm base})(T)
   =\rho_b\ell(F'_b/2)\int_1^T K_T(T,x)m(x)\,\dd x,\nonumber\\
 K(T,x)&=\frac{\sinh(2\ell(T-x))}{2\ell}
                       \frac{\cosh^2(\ell x)}{\cosh^2(\ell T)},\qquad
 K_T=\frac{\cosh^2(\ell x)\cosh(\ell T-2\ell x)}{\cosh^3(\ell T)}.
 \label{eq:kernels}
\end{align}
The exact bounds are $0\le K\le T-x$ and $0<K_T\le1$.
The factor $1/2$ in these kernels' source coefficient is retained; there
is no extra factor multiplying the differences of the observable second
derivatives.

\begin{theorem}[Conditional integrated geometric control]\label{thm:main}
Under the stated retained premises, each admitted smooth compact
preparation satisfies
\begin{equation}
 |\nu(T)|\le C_vR,\qquad |\vartheta(T)|\le C_\Theta R,\qquad
 C_v=\frac{17161}{19200000000000000000},\quad
 C_\Theta=\frac{131}{2400000000000}.
 \label{eq:transfer}
\end{equation}
For the specified closed class
\begin{equation}
 \mathcal C_{H^1}=\{h\text{ satisfying \eqref{eq:admissible}}:
                         \norm D_{H^1(I)}\le2\times10^{31}\},
 \label{eq:class}
\end{equation}
the coefficients at $\tau_*$ obey the outward enclosures
\begin{equation}
 v_2\in[-7.2400\times10^{16},-8.6400\times10^{15}],\qquad
 \Theta_2\in[-1.0113\times10^{22},-3.3540\times10^{21}].
 \label{eq:signs}
\end{equation}
For every individual admitted exact family, for sufficiently small
nonzero $\eps$ of either sign,
\begin{equation}
 0<V_\eps(\tau_*)<V_0(\tau_*),\qquad
 0<\Theta_\eps(\tau_*)<\Theta_0(\tau_*).
 \label{eq:nonlinear}
\end{equation}
The sufficiently small interval can depend on the family.
\end{theorem}

\begin{proof}
Equations \eqref{eq:trace} and \eqref{eq:kernels} imply
\begin{equation}
 C_v=\frac{(131\ell)^2(1/1000)\Delta^2}{4\cdot300},\qquad
 C_\Theta=\frac{131\ell(1/1000)\Delta}{2\cdot300},
 \label{eq:Cformula}
\end{equation}
which equal \eqref{eq:transfer}. At $R=2\times10^{31}$ the absolute
increment caps are below $1.78761\times10^{16}$ and
$1.09167\times10^{21}$. Adding and subtracting the exact caps from the
retained base rational intervals, whose outward displays are
\begin{equation}
 v_2^{\rm base}\in[-5.4520,-2.6521]\,10^{16},\qquad
 \Theta_2^{\rm base}\in[-9.0208,-4.4458]\,10^{21},
 \label{eq:baseintervals}
\end{equation}
gives \eqref{eq:signs}. The radius and energy/trace gates were declared
before evaluation of these final sign gates.

The first integrated volume and expansion variations vanish by the shared
quadrupole mean. For each fixed family, retained $C^2$ dependence gives
\begin{equation}
 V_\eps/V_0=1+\tfrac12\eps^2v_2+o(\eps^2),\qquad
 \Theta_\eps=\Theta_0+\tfrac12\eps^2\Theta_2+o(\eps^2).
 \label{eq:Taylor}
\end{equation}
Strict negativity and continuity, with $V_0,\Theta_0>0$ at the target,
give \eqref{eq:nonlinear} after a family-dependent shrinkage.
\end{proof}

The estimating norm controls no higher derivatives or nonlinear remainders
uniformly. The theorem therefore gives no effective or class-wide nonlinear
parameter interval, signed time history, contracting branch, or
volume-acceleration sign. Large second coefficients are not finite-amplitude
field values.

\section{Actual preparation data beyond the old supremum criterion}

The earlier sufficient condition used
$R_A=\max\{\sup_I|D|,\sup_I|D_r|\}\le2\times10^{33}$
and an initial supremum quantity $N$ involving $U_{0,r}$
\cite{intrinsic}. A direct $H^1$ criterion cannot be obtained merely by
embedding into that condition. This failure occurs in the actual retained
preparation image, not just among arbitrary scalar profiles.

For any smooth $Z$ compactly supported inside $I$, define
$\chi=\ell\rho^3Z$ and $h=3bZ-Z_r$. Then
$\chi_r=-\ell\rho^3h$ and $\chi_y=\rho^2h$, so the exact moment in
\eqref{eq:admissible} vanishes by endpoint evaluation. With
\begin{equation}
 \mathfrak A=\frac3{2g},\qquad c=-2b-g,\qquad
 d=3(\ell^2-b^2+bg-b_r),
 \label{eq:profileoperator}
\end{equation}
elimination of $h$ in \eqref{eq:reconstruct} gives
\begin{equation}
 D=\mathfrak A(Z_{rr}+cZ_r+dZ),\quad
 D_r=\mathfrak A Z_{rrr}+(\mathfrak A_r+\mathfrak Ac)Z_{rr}
       +(\mathfrak A_rc+\mathfrak Ac_r+\mathfrak Ad)Z_r
       +(\mathfrak A_rd+\mathfrak Ad_r)Z.
 \label{eq:Doperator}
\end{equation}
The coefficients are bounded on a compact interior interval using
$b\in C^2$, $g\in C^1$ and $g>0$. Literal smooth admissibility uses the
retained smooth static background, distinct from the lower $C^4$
topology for parameter dependence.

Take $r_0=9/8$, integers $m\ge4$, and
\begin{equation}
 \psi(u)=
 \begin{cases}u^3\exp(1-1/(1-u^2)),&|u|<1,\\0,&|u|\ge1,\end{cases}
 \qquad Z_m(r)=m^{-5}\psi(m^2(r-r_0)).
 \label{eq:concentration}
\end{equation}
This is smooth and flat, supported inside
$[17/16,19/16]\subset(1,T)$. The raw metric addition and conditioning
generator are consequently compact and flat, with the other canonical
sectors still fixed. Change of variables gives the exact scaling
\begin{equation}
 \norm{Z_m^{(j)}}_{L^2(I)}
       =m^{2j-6}\norm{\psi^{(j)}}_{L^2(-1,1)},\qquad 0\le j\le3.
 \label{eq:scaling}
\end{equation}
Therefore \eqref{eq:Doperator} bounds $\norm{D_m}_{H^1(I)}$ by one
finite constant $C_*>0$ independent of $m$.
At the center,
$\psi(0)=\psi'(0)=\psi''(0)=0$ and $\psi'''(0)=6$, so exactly
\begin{equation}
 h_m(r_0)=h_{m,r}(r_0)=D_m(r_0)=0,\qquad
 D_{m,r}(r_0)=\frac{9m}{g(r_0)},\qquad
 w_{m,r}(r_0)=-\frac{3p(r_0)m}{a(r_0)g(r_0)}.
 \label{eq:center}
\end{equation}
Thus $R_A$ and the old genuine initial derivative supremum, hence every
sufficient old $N$ of that form, diverge. The profile amplitudes themselves
can tend to zero; the failure concerns derivative suprema.

For any prescribed positive budget $B$, multiplying every direction by
the same $\eta=B/(2C_*)>0$ puts its $H^1$ norm below $B$ and preserves
this divergence. With $B=2\times10^{31}$, infinitely many such data
satisfy Theorem~\ref{thm:main} while failing the old $R_A$ cutoff.
This proves coverage beyond the previous sufficient criterion. It does
not state that the new budget ball contains the old one, that the actual
responses diverge, or that sign instability has been found. Each fixed
$m$ has only its retained family-dependent nonlinear interval.

\section{Preparation selection and relation to established methods}

The energy proof is a tailored application of standard multiplier methods.
Kreiss, Reula, Sarbach and Winicour construct positive mixed-derivative
energies, boundary dissipation and lower-order estimates in
\cite[Sec.~2 and Apps.~1--2]{KRSW2007}. Their geometric treatment
\cite{KRSW2009} likewise adapts the timelike multiplier to the boundary.
Their theorems impose their own boundary conditions; the 2007 displayed
condition has a strictly positive time-derivative coefficient.
Neither is invoked as a theorem automatically covering the present
mixed Dirichlet/nonselfadjoint Robin system or shrinking cone.
Equations \eqref{eq:flux}--\eqref{eq:E0} provide that direct conditional
estimate with the retained model-specific caps.

The estimating energy also does not select the preparation budget.
A specified first-tangent quadratic functional $Q(v)$ is common to
every preparation here because $v$ is common. Hollands--Wald canonical
energy explicitly has such first-tangent dependence in its stationary
black-hole setting \cite[Sec.~2.3]{HW}; that setting and its Killing
generator are not supplied by this shell's proper clock.
The retained matched-shell Noether calculation supplies necessary
quadratic tests of the form
\begin{equation}
 DH_K|_0=0,\qquad D^2H_K|_0[v,v]
       =\Omega_{{\rm rel},0}(v,\mathcal K v)=0
 \label{eq:Noether}
\end{equation}
under its own full-domain and compensated-symmetry premises
\cite{Noether}. Such tests are constant on the admitted second-preparation
fiber. They cannot alone select a finite $H^1$ radius: scaling a fixed
nonzero interior profile by any fixed finite parameter changes its $R$
without changing these first-tangent inputs.

The limitation is specific. For an arbitrary $C^2$ nonlinear functional,
\begin{equation}
 (F\circ x_\eps)''|_0=DF_0[w]+D^2F_0[v,v].
 \label{eq:functional}
\end{equation}
The first term can distinguish preparations. Nor is a bare Noether
surface integral automatically a Hamiltonian: generator, boundary
conditions, shell and corner completion, and allowed variations matter
\cite{IW,HarlowWu}. A specified local initial-cut functional depending
on identical finite boundary and embedding jets has equal coefficients
through order two by the chain rule, since these compact additions
vanish on the initial collars. This does not apply to arbitrary
interior or future cuts, interior-dependent generators, or unspecified
physical energies. No no-go theorem for every energy or state-based
preparation rule is claimed.

The physical initial section, continued component, scalar-label
calibration, canonical-sector restriction and numerical radius remain
inputs. The present result certifies what follows from one integrated
geometric specification; it supplies no action, state or process that
chooses it. No measured amplitude, matter/radiation energy transfer or
higher-dimensional Big Bang origin is established. The focused primary
comparison identifies the standard methods used here, without asserting
exhaustive novelty or priority.

\appendix
\section{Certificates and reproducibility}

The companion capsule \cite{records} contains the tilted-energy proof,
predeclared design, selected-source SHA-256 pins, exact rational producer,
normal and optimized reports, independent reconstruction, actual-image
concentration proof and its algebra checker, and the limited
preparation-selection and primary-source comparisons.
The producer authenticates retained records before checking rational
consequences. Authentication of a published coefficient assertion is not
a fresh proof of that continuous assertion.

The 24 finite arithmetic gates verify the positive tilt, bulk and nonsymmetric
Robin rate guards, exponential remainder, complete initial map, boundary
trace factor, and whole-class sign displays. The concentration checker
verifies operator identities, center jets, support exponents and
continuous analytic bump bounds without evaluating a numerical bump.
No new background/evolved-field oracle, PDE solve or numerical bump query
is used. Independent implementations share the retained analytic and
background premises; their agreement is not an independent observation
of the nonlinear PDE. Source, displayed-bound audit, deterministic
compilation receipt and page-render review accompany this working draft.

\begin{thebibliography}{10}
\footnotesize
\setlength{\itemsep}{2pt}
\bibitem{KRSW2007} H.-O. Kreiss, O. Reula, O. Sarbach and J. Winicour,
\emph{Well-posed initial-boundary value problem for the harmonic Einstein
equations using energy estimates}, Class. Quantum Grav. \textbf{24}
(2007) 5973--5984,
\href{https://arxiv.org/abs/0707.4188}{arXiv:0707.4188}.
\bibitem{KRSW2009} H.-O. Kreiss, O. Reula, O. Sarbach and J. Winicour,
\emph{Boundary conditions for coupled quasilinear wave equations with
application to isolated systems}, Commun. Math. Phys. \textbf{289}
(2009) 1099--1129,
\href{https://arxiv.org/abs/0807.3207}{arXiv:0807.3207}.
\bibitem{HW} S. Hollands and R. M. Wald,
\emph{Stability of Black Holes and Black Branes}, Commun. Math. Phys.
\textbf{321} (2013) 629--680,
\href{https://arxiv.org/abs/1201.0463}{arXiv:1201.0463}.
\bibitem{IW} V. Iyer and R. M. Wald,
\emph{Some Properties of Noether Charge and a Proposal for Dynamical Black
Hole Entropy}, Phys. Rev. D \textbf{50} (1994) 846--864,
\href{https://arxiv.org/abs/gr-qc/9403028}{arXiv:gr-qc/9403028}.
\bibitem{HarlowWu} D. Harlow and J.-q. Wu,
\emph{Covariant phase space with boundaries}, JHEP \textbf{10} (2020) 146,
\href{https://arxiv.org/abs/1906.08616}{arXiv:1906.08616}.
\bibitem{preparation} R. Maldonado, \emph{Preparation dependence of
proper-time volume in Einstein--scalar shell perturbations}, working
manuscript (4 November 2026),
\href{https://github.com/maldonado-research/HDblast/tree/main/research/HDBLAST_SHELL_PREPARATION_DEPENDENCE_20261104}{preparation-difference capsule}.
\bibitem{intrinsic} R. Maldonado, \emph{Scalar-level initial geometry and
robust shell-volume perturbations}, working manuscript (4 November 2026),
\href{https://github.com/maldonado-research/HDblast/tree/main/research/HDBLAST_INTRINSIC_PREPARATION_ROBUSTNESS_20261104}{scalar-level preparation capsule}.
\bibitem{Noether} R. Maldonado,
\emph{Full nonlinear shell Noether identity and quadratic integrability
tests}, conditional technical record (13 October 2026),
\href{https://github.com/maldonado-research/HDblast/tree/main/research/HDBLAST_SHELL_NOETHER_20261013}{shell-Noether capsule}.
\bibitem{records} R. Maldonado, \emph{HDBLAST conditional technical
records and reproducibility materials} (2026),
\href{https://github.com/maldonado-research/HDblast}{repository}.
The retained source of the operator assertions is the
23 October total-response-preflight capsule.
The present companion is
\href{https://github.com/maldonado-research/HDblast/tree/main/research/HDBLAST_H1_PREPARATION_RESPONSE_20261104}{the integrated geometric response capsule}.
\end{thebibliography}
\end{document}
